\documentclass[11pt]{article}

\usepackage[margin=1in]{geometry}
\usepackage{booktabs}
\usepackage{longtable}
\usepackage{array}
\usepackage{amsmath}
\usepackage{graphicx}
\usepackage{pdflscape}
\usepackage{hyperref}
\usepackage{xcolor}

\hypersetup{
  colorlinks=true,
  linkcolor=blue,
  urlcolor=blue
}

\title{Joint versus Separate Optimisation: 100-Replication Simulation Review}
\author{gamlss.longitudinal JSS replication notes}
\date{Compiled from outputs completed 2026-06-14}

\begin{document}
\maketitle

\section{Purpose}

This note summarises the 100-replication simulation comparing the two
random-search optimisation strategies used in the JSS replication workflow:
\texttt{rs\_joint} and \texttt{rs\_separate}.  The goal was to test whether
joint optimisation gives a practically meaningful improvement over separate
optimisation once the iteration caps are increased enough to give the joint
method a fairer convergence opportunity.

\section{Run Provenance}

The completed checkpointed run wrote the final combined outputs on
\textbf{2026-06-14 22:49:15 +1000}.  The checkpointing design wrote one CSV
per replication before combining results, so machine restarts no longer discard
all in-memory progress.

\begin{table}[ht]
\centering
\caption{Completed run artifacts.}
\begin{tabular}{ll}
\toprule
Item & Value \\
\midrule
Replications & 100 \\
Rows per replication & 36 \\
Total result rows & 3600 \\
\texttt{rs\_joint} fits & 1800 \\
\texttt{rs\_separate} fits & 1800 \\
Checkpoint directory &
\texttt{results/jss-replication/expanded/checkpoints/03-joint-vs-separate-optimization-100rep-outer16} \\
Final result CSV &
\texttt{results/jss-replication/expanded/data/03-joint-vs-separate-optimization-results.csv} \\
Summary table &
\texttt{results/jss-replication/expanded/tables/03-joint-vs-separate-optimization-summary.csv} \\
Metric-win table &
\texttt{results/jss-replication/expanded/tables/03-joint-vs-separate-optimization-metric-wins.csv} \\
\bottomrule
\end{tabular}
\end{table}

\section{Simulation Design}

The formal run used the expanded profile for module 03,
\texttt{03-joint-vs-separate-optimization}.  The simulation grid was built
with \texttt{run\_coverage\_simulations()} and then fit with only the two
longitudinal random-search methods, \texttt{rs\_separate} and
\texttt{rs\_joint}.  Table~\ref{tab:settings} lists the settings used in the
completed run.

\begin{table}[ht]
\centering
\caption{Simulation settings used in the 100-replication run.}
\label{tab:settings}
\begin{tabular}{>{\ttfamily}ll}
\toprule
Setting & Value and role \\
\midrule
profile & \texttt{expanded}; writes under \texttt{results/jss-replication/expanded}. \\
families & \texttt{NO}, \texttt{NBI}. \texttt{NO} is the normal response family; \texttt{NBI} is the negative-binomial type I count family. \\
copulas & \texttt{N}, \texttt{C}, \texttt{G}. These are Gaussian, Clayton, and Gumbel copulas. \\
designs & \texttt{intercept}, \texttt{covariate}, \texttt{time\_dependence}. Details are given below. \\
n & 48 subjects per simulated data set. \\
times & Four repeated measurements, \(\{1,2,3,4\}\), per subject before missingness. \\
dependence & \texttt{strong}; used for constant-dependence designs. \\
missingness & \texttt{none}; no responses were removed after simulation. \\
methods & \texttt{rs\_separate}, \texttt{rs\_joint}. \\
start\_mode & \texttt{default}; no truth-adjacent warm start was requested. \\
max\_outer\_iter & 16. \\
max\_inner\_iter & 8. \\
max\_elapsed\_sec & 120 seconds per fit attempt. \\
base seed & 20260528. \\
\bottomrule
\end{tabular}
\end{table}

\subsection{Response Families}

The response distribution was generated from \texttt{gamlss.dist} families.
For the two families used here, the coverage simulator applies conservative
default parameters:

\begin{itemize}
  \item \texttt{NO}: \(\mu = 0\), \(\sigma = 1\).
  \item \texttt{NBI}: \(\mu = 4\), \(\sigma = 0.5\).
\end{itemize}

These base values are used directly in the \texttt{intercept} and
\texttt{time\_dependence} designs.  In the \texttt{covariate} design, the
\(\mu\) parameter is modified by a subject-level covariate as described below.

\subsection{Copula Families and Dependence Strength}

The simulator uses a first-order longitudinal copula model between adjacent
time points.  The included copula codes were:

\begin{itemize}
  \item \texttt{N}: Gaussian copula, parameterised by a correlation-like
  \(\theta\) with Fisher-\(z\) link in the fitted model.
  \item \texttt{C}: Clayton copula, parameterised internally through Kendall's
  \(\tau\) in the simulator and a positive \(\theta\) link in the fitted model.
  \item \texttt{G}: Gumbel copula, also generated from a Kendall's \(\tau\)
  path and fit through the Gumbel parameter link.
\end{itemize}

For constant-dependence designs, the \texttt{strong} dependence setting maps to
\(\rho = 0.65\) for the Gaussian copula and \(\tau = 0.55\) for the Clayton and
Gumbel copulas.  For the \texttt{time\_dependence} design, the constant
\texttt{strong} setting is replaced by a deterministic time-varying dependence
path:

\begin{itemize}
  \item \texttt{N}: \(\theta(t) = 0.1 + 0.45\,s(t)\), where \(s(t)\) rescales
  the left endpoint of each adjacent time pair to \([0,1]\).
  \item \texttt{C} and \texttt{G}: \(\tau(t) = 0.1 + 0.3\,s(t)\).
\end{itemize}

With four time points, the adjacent pair left endpoints are \(1,2,3\), so the
time-varying path is evaluated over three adjacent-pair dependence levels.

\subsection{Designs}

The three design settings control which part of the data-generating model is
non-constant and which formulas are used when fitting each method.

\paragraph{\texttt{intercept}.}
All marginal parameters and the copula parameter are constant.  The fitted
formulas are intercept-only for the margins and dependence:
\[
\mu \sim 1,\quad \sigma \sim 1,\quad \nu \sim 1,\quad \tau \sim 1,\quad
\theta \sim 1,\quad \zeta \sim 1.
\]

\paragraph{\texttt{covariate}.}
A subject-level covariate \(x\) is generated by rescaling the subject index to
\([0,1]\) and then centering it.  The data-generating \(\mu\) parameter varies
with this covariate.  For positive baseline means, the simulator uses
\[
\mu(x) = \mu_0 \exp(0.15x),
\]
while for non-positive baseline means it uses
\[
\mu(x) = \mu_0 + 0.15x.
\]
Thus, for \texttt{NO}, \(\mu(x)=0.15x\); for \texttt{NBI},
\(\mu(x)=4\exp(0.15x)\).  The fitted formulas use \(x\) only in the mean
submodel:
\[
\mu \sim x,\quad \sigma \sim 1,\quad \nu \sim 1,\quad \tau \sim 1,\quad
\theta \sim 1,\quad \zeta \sim 1.
\]

\paragraph{\texttt{time\_dependence}.}
The marginal response model remains intercept-only, but the adjacent-time
copula dependence changes across time.  The fitted dependence formula is
\[
\theta \sim \text{time},
\]
with the other marginal parameters held at intercept-only formulas.  This
design is the one that directly tests recovery of time-varying dependence; it
is also why \texttt{benchmark\_theta\_time\_abs\_error} is finite only for the
time-dependence scenarios.

\subsection{Replication and Seeding}

For replication \(r\), the module called
\texttt{run\_coverage\_simulations()} with seed
\[
20260528 + 3000 + 100r.
\]
Inside the coverage grid runner, the individual case seed is offset again by
the case index \(i\), giving effective case-level seeds
\[
20260528 + 3000 + 100r + i.
\]
This made each family--copula--design--method cell reproducible while avoiding
identical random streams across cells.

\subsection{Optimisation Settings}

Both methods were run with the same outer, inner, and elapsed-time caps.  The
\texttt{rs\_separate} method optimises marginal and dependence components in a
separate/backfitting-style sequence, while \texttt{rs\_joint} evaluates the
joint marginal-plus-copula objective in the random-search loop.  The increased
\texttt{max\_outer\_iter = 16} and \texttt{max\_elapsed\_sec = 120} settings
were chosen after the 10-replication calibration showed that the lower-cap
joint run had avoidable non-convergence.

Each replication therefore contained
\[
2 \text{ families} \times 3 \text{ copulas} \times 3 \text{ designs}
\times 2 \text{ methods} = 36
\]
fits.  Across 100 replications, this gives 1800 paired comparison cells and
3600 fitted-model rows.

\section{Candidate Set}

The module first mined retained exploratory summaries and selected seven
candidate cases for review.  These candidates determined the copula set used by
the formal grid.  Five were classified as \texttt{joint\_win}, one as
\texttt{tie\_control}, and one as \texttt{cautionary}.

\begin{table}[ht]
\centering
\caption{Selected candidate sources.}
\begin{tabular}{llll}
\toprule
Candidate & Role & Copula & Source \\
\midrule
\texttt{jvs-01} & \texttt{joint\_win} & \texttt{N} &
\texttt{top\_cases\_loglik\_logscore\_all\_rmse\_improved} \\
\texttt{jvs-02} & \texttt{joint\_win} & \texttt{N} &
\texttt{stress\_balanced\_candidates} \\
\texttt{jvs-03} & \texttt{joint\_win} & \texttt{C} &
\texttt{top\_cases\_loglik\_logscore\_all\_rmse\_improved} \\
\texttt{jvs-04} & \texttt{joint\_win} & \texttt{N} &
\texttt{stress\_balanced\_candidates} \\
\texttt{jvs-05} & \texttt{joint\_win} & \texttt{N} &
\texttt{stress\_balanced\_candidates} \\
\texttt{jvs-06} & \texttt{tie\_control} & \texttt{G} &
\texttt{top\_cases\_loglik\_logscore\_all\_rmse\_improved} \\
\texttt{jvs-07} & \texttt{cautionary} & \texttt{C} &
\texttt{joint\_vs\_separate\_delta\_summary} \\
\bottomrule
\end{tabular}
\end{table}

\section{Overall Completion, Convergence, and Runtime}

Both methods completed most fits, but the separate optimiser retained a clear
convergence advantage.  The joint optimiser's convergence improved relative to
the earlier low-cap pilot, but did not reach parity with the separate method.

\begin{table}[ht]
\centering
\caption{Overall method-level completion and runtime.}
\begin{tabular}{lrrrrr}
\toprule
Method & Fits & Success rate & Convergence rate & Median sec. & Mean sec. \\
\midrule
\texttt{rs\_joint} & 1800 & 0.966 & 0.861 & 5.10 & 15.68 \\
\texttt{rs\_separate} & 1800 & 0.983 & 0.982 & 1.71 & 9.35 \\
\bottomrule
\end{tabular}
\end{table}

\begin{table}[ht]
\centering
\caption{Completion and runtime by response family.}
\begin{tabular}{llrrrr}
\toprule
Method & Family & Success rate & Convergence rate & Median sec. & Mean sec. \\
\midrule
\texttt{rs\_joint} & \texttt{NBI} & 0.932 & 0.863 & 13.14 & 28.17 \\
\texttt{rs\_separate} & \texttt{NBI} & 0.967 & 0.967 & 5.67 & 17.68 \\
\texttt{rs\_joint} & \texttt{NO} & 1.000 & 0.858 & 2.31 & 3.19 \\
\texttt{rs\_separate} & \texttt{NO} & 1.000 & 0.998 & 0.74 & 1.02 \\
\bottomrule
\end{tabular}
\end{table}

\subsection{Additional Fit Outcomes}

The result table also records diagnostic fit outcomes beyond success,
convergence, and runtime.  These include convergence-timeout classifications,
warnings, marginal fallback use, likelihood components, parameter recovery
errors, and fitted dependence summaries.  Table~\ref{tab:additional-fit}
summarises the most informative populated fields.

\begin{table}[ht]
\centering
\caption{Additional method-level fit outcomes.}
\label{tab:additional-fit}
\begin{tabular}{lrr}
\toprule
Outcome & \texttt{rs\_joint} & \texttt{rs\_separate} \\
\midrule
Convergence-timeout rate & 0.034 & 0.017 \\
Warning rate & 0.188 & 0.083 \\
Marginal fallback rate & 0.000 & 0.000 \\
Median joint log-likelihood & -269.50 & -273.19 \\
Median marginal log-likelihood & -297.26 & -299.36 \\
Median copula log-likelihood & 41.63 & 40.29 \\
Median max. absolute parameter error & 0.333 & 0.336 \\
Median absolute copula-\(\tau\) error & 0.055 & 0.057 \\
90th percentile elapsed sec. & 42.14 & 23.44 \\
\bottomrule
\end{tabular}
\end{table}

All recorded failures were classified as convergence timeouts; there were no
marginal fallback fits in either method, and all fits used the default attempt
path.  The unpaired method-level medians suggest slightly higher likelihoods
and slightly smaller parameter errors for \texttt{rs\_joint}, but the paired
finite deltas were close to neutral: mean paired joint-log-likelihood delta was
\(-0.32\), median paired joint-log-likelihood delta was \(-0.35\), and joint
won 49.1\% of finite paired likelihood comparisons.  For maximum absolute
parameter error, joint had a small paired advantage, with mean separate-minus-
joint error \(0.0031\), median \(0.0029\), and a 51.0\% joint win rate.  For
absolute copula-\(\tau\) error, the sign reversed slightly: mean
separate-minus-joint error was \(-0.0021\), median \(-0.0012\), and the joint
win rate was 48.9\%.

\section{Metric-Win Results}

The metric-win table compares methods within paired simulation cells.  For
error and scoring metrics, lower raw values are better; for
\texttt{benchmark\_pit\_ks\_p\_value}, larger values are better.  Runtime is
also treated as a metric where lower is better.  Some metrics have fewer finite
comparisons because failures or non-finite benchmark quantities remove paired
cells from that metric.

\begin{table}[ht]
\centering
\caption{Selected metric win and win-or-tie rates.}
\begin{tabular}{llrrrr}
\toprule
Metric & Method & Finite \(n\) & Wins & Win rate & Win/tie rate \\
\midrule
\texttt{neg\_log\_score} & \texttt{rs\_joint} & 1739 & 871 & 0.492 & 0.785 \\
\texttt{neg\_log\_score} & \texttt{rs\_separate} & 1770 & 901 & 0.508 & 0.796 \\
\texttt{mean\_rmse} & \texttt{rs\_joint} & 1739 & 879 & 0.496 & 0.517 \\
\texttt{mean\_rmse} & \texttt{rs\_separate} & 1770 & 893 & 0.504 & 0.517 \\
\texttt{q90\_mae} & \texttt{rs\_joint} & 1739 & 1032 & 0.582 & 0.597 \\
\texttt{q90\_mae} & \texttt{rs\_separate} & 1770 & 1021 & 0.576 & 0.589 \\
\texttt{theta\_time\_abs\_error} & \texttt{rs\_joint} & 600 & 309 & 0.515 & 0.530 \\
\texttt{theta\_time\_abs\_error} & \texttt{rs\_separate} & 600 & 291 & 0.485 & 0.495 \\
\texttt{lower\_tail\_error\_05} & \texttt{rs\_joint} & 1739 & 1309 & 0.739 & 0.739 \\
\texttt{lower\_tail\_error\_05} & \texttt{rs\_separate} & 1770 & 1387 & 0.783 & 0.783 \\
\texttt{upper\_tail\_error\_05} & \texttt{rs\_joint} & 1739 & 939 & 0.530 & 0.532 \\
\texttt{upper\_tail\_error\_05} & \texttt{rs\_separate} & 1770 & 1027 & 0.580 & 0.580 \\
\texttt{pit\_mean\_abs\_error} & \texttt{rs\_joint} & 1739 & 792 & 0.447 & 0.507 \\
\texttt{pit\_mean\_abs\_error} & \texttt{rs\_separate} & 1770 & 980 & 0.553 & 0.601 \\
\texttt{interval\_coverage\_95} & \texttt{rs\_joint} & 1739 & 960 & 0.542 & 0.542 \\
\texttt{interval\_coverage\_95} & \texttt{rs\_separate} & 1770 & 1089 & 0.615 & 0.615 \\
\texttt{elapsed\_sec} & \texttt{rs\_joint} & 1800 & 117 & 0.065 & 0.081 \\
\texttt{elapsed\_sec} & \texttt{rs\_separate} & 1800 & 1683 & 0.935 & 0.943 \\
\bottomrule
\end{tabular}
\end{table}

\section{Scenario-Level Deltas}

Table~\ref{tab:scenario-deltas} reports scenario-level averages for the paired
comparison.  The delta columns use the sign convention:

\[
\Delta_{\mathrm{error}} =
\text{separate error} - \text{joint error},
\]
so positive values favour the joint method for error and score metrics.  For
joint log-likelihood,
\[
\Delta_{\log L} =
\log L_{\mathrm{joint}} - \log L_{\mathrm{separate}},
\]
so positive values also favour the joint method.  For elapsed time,
\[
\Delta_{\mathrm{time}} =
\text{separate sec.} - \text{joint sec.},
\]
so negative values indicate that the joint method was slower.

\begin{landscape}
\begin{longtable}{lllrrrrrrrr}
\caption{Scenario-level paired deltas, ordered by core win rate.}
\label{tab:scenario-deltas}\\
\toprule
Family & Copula & Design & Joint conv. & Sep. conv. & Core win &
\(\Delta\) NLS & \(\Delta\) mean RMSE & \(\Delta\) q90 &
\(\Delta\theta(t)\) & \(\Delta\) sec. \\
\midrule
\endfirsthead
\toprule
Family & Copula & Design & Joint conv. & Sep. conv. & Core win &
\(\Delta\) NLS & \(\Delta\) mean RMSE & \(\Delta\) q90 &
\(\Delta\theta(t)\) & \(\Delta\) sec. \\
\midrule
\endhead
\texttt{NO} & \texttt{G} & \texttt{intercept} & 0.74 & 1.00 & 0.533 & -0.011 & 0.018 & 0.021 & -- & -3.219 \\
\texttt{NO} & \texttt{N} & \texttt{intercept} & 1.00 & 1.00 & 0.525 & 0.013 & 0.019 & 0.011 & -- & -1.476 \\
\texttt{NO} & \texttt{C} & \texttt{time\_dependence} & 1.00 & 0.99 & 0.508 & 0.000 & 0.001 & -0.001 & 0.040 & -1.396 \\
\texttt{NO} & \texttt{G} & \texttt{time\_dependence} & 1.00 & 0.99 & 0.504 & -0.007 & 0.012 & 0.003 & -0.019 & -1.307 \\
\texttt{NO} & \texttt{C} & \texttt{intercept} & 0.58 & 1.00 & 0.498 & 0.010 & -0.013 & 0.016 & -- & -2.973 \\
\texttt{NO} & \texttt{N} & \texttt{covariate} & 1.00 & 1.00 & 0.495 & -0.017 & 0.010 & 0.007 & -- & -1.452 \\
\texttt{NO} & \texttt{G} & \texttt{covariate} & 0.77 & 1.00 & 0.478 & 0.004 & 0.012 & 0.006 & -- & -3.263 \\
\texttt{NO} & \texttt{N} & \texttt{time\_dependence} & 1.00 & 1.00 & 0.464 & 0.011 & -0.003 & -0.007 & -0.006 & -0.674 \\
\texttt{NO} & \texttt{C} & \texttt{covariate} & 0.63 & 1.00 & 0.455 & -0.012 & -0.014 & -0.003 & -- & -3.751 \\
\texttt{NBI} & \texttt{G} & \texttt{time\_dependence} & 0.98 & 0.98 & 0.400 & 0.034 & 0.042 & 0.020 & 0.034 & -2.776 \\
\texttt{NBI} & \texttt{N} & \texttt{covariate} & 0.43 & 0.73 & 0.378 & -0.017 & 0.024 & 0.062 & -- & -23.441 \\
\texttt{NBI} & \texttt{N} & \texttt{time\_dependence} & 1.00 & 1.00 & 0.362 & 0.005 & 0.027 & -0.060 & 0.013 & -18.264 \\
\texttt{NBI} & \texttt{C} & \texttt{intercept} & 0.89 & 1.00 & 0.360 & -0.004 & 0.044 & 0.180 & -- & -5.776 \\
\texttt{NBI} & \texttt{G} & \texttt{intercept} & 0.85 & 1.00 & 0.355 & -0.004 & 0.095 & 0.180 & -- & -6.504 \\
\texttt{NBI} & \texttt{C} & \texttt{time\_dependence} & 0.99 & 0.99 & 0.352 & -0.007 & 0.014 & 0.160 & -0.058 & -2.882 \\
\texttt{NBI} & \texttt{N} & \texttt{intercept} & 1.00 & 1.00 & 0.348 & 0.003 & 0.025 & 0.090 & -- & -17.566 \\
\texttt{NBI} & \texttt{G} & \texttt{covariate} & 0.76 & 1.00 & 0.313 & -0.010 & -0.140 & -0.153 & -- & -9.353 \\
\texttt{NBI} & \texttt{C} & \texttt{covariate} & 0.87 & 1.00 & 0.295 & -0.012 & -0.089 & -0.118 & -- & -7.847 \\
\bottomrule
\end{longtable}
\end{landscape}

\section{Interpretation}

The 100-replication run gives a more conservative picture than the earlier
10-replication calibration.  The joint method is competitive on selected
prediction and dependence-recovery metrics, but it does not dominate the
separate method overall.

\paragraph{Where joint optimisation helped most.}
The clearest joint-favouring scenarios were in the Gaussian/normal family.  The
\texttt{NO + G + intercept} and \texttt{NO + N + intercept} scenarios had the
highest core win rates.  The \texttt{NO + C + time\_dependence} case is still a
useful example because both methods converged reliably and the joint method
showed a positive mean gain on the time-varying dependence error.  Across all
finite paired cells, the joint method had a small edge on
\texttt{benchmark\_q90\_mae} and \texttt{benchmark\_theta\_time\_abs\_error}.

\paragraph{Where separate optimisation remained stronger.}
The separate method remained substantially faster.  It won or tied on runtime
in 94.3\% of paired comparisons, while the joint method won or tied in only
8.1\%.  Separate optimisation also retained better global convergence:
98.2\% versus 86.1\%.  Calibration and tail metrics generally leaned separate,
especially \texttt{benchmark\_pit\_mean\_abs\_error},
\texttt{benchmark\_tail\_error\_lower\_05}, and
\texttt{benchmark\_tail\_error\_upper\_05}.

\paragraph{Family-level pattern.}
The \texttt{NBI} scenarios were materially harder.  Joint optimisation had
lower success and convergence rates for \texttt{NBI}, and some
\texttt{NBI + N} cells were especially slow.  Even where the joint method
improved some predictive metrics in \texttt{NBI}, it often did so with much
larger runtime costs and weaker convergence.

\paragraph{Bottom line.}
The full run supports a nuanced conclusion: increasing the iteration cap made
joint optimisation viable enough for review, but it did not make joint
optimisation the uniformly preferred method.  Joint optimisation is best framed
as a targeted option that can improve upper-quantile prediction or
time-varying dependence recovery in some Gaussian scenarios, while separate
optimisation remains the more stable and computationally efficient default.

\section{Figures}

The completed run also produced two diagnostic figures:

\begin{itemize}
  \item
  \texttt{results/jss-replication/expanded/figures/03-joint-vs-separate-optimization-deltas.png}
  \item
  \texttt{results/jss-replication/expanded/figures/03-joint-vs-separate-optimization-metric-dashboard.png}
\end{itemize}

These are not embedded here to keep the LaTeX source portable across working
directories, but they should be reviewed alongside this note.

\end{document}
